% GENERATED FILE. Edit canonical lessons, then run npm run generate:books.
% canonical-source-hash: fnv1a64:8bd695245e33ba77
% canonical-lessons: AR-C548-sibaha, AR-C548-mulakama, AR-C548-saraa, AR-C548-nafasa, AR-C548-lakama

\chapter{In the pool and the ring}
\label{ch:ar-in-the-pool-and-the-ring}
\hlchaptermodality{548}

\begin{chapteropening}
\textbf{By the end of this chapter:} \emph{I can talk about swimming and boxing, and say that someone wrestled, competed with others or punched.}

Complete the last lesson of chapter 548: I can talk about swimming and boxing, and say that someone wrestled, competed with others or punched.
\end{chapteropening}

\section[\texorpdfstring{\ar{سباحة} (sibāḥa)}{سباحة (sibāḥa)}]{\ar{سباحة} (sibāḥa) — swimming}

\label{lesson:AR-C548-sibaha}



\begin{quote}
Before the new one: say the Arabic for a medal, then the Arabic for to qualify.
\end{quote}

\subsection*{You'll want to know: \ar{سباحة}}
\textbf{\ar{سباحة}} — \emph{sibāḥa} — \textquotedblleft{}swimming\textquotedblright{}.

\begin{quote}
\ar{السباحة} \ar{رياضة} \ar{صحية}.
\end{quote}

\begin{quote}
\emph{as-sibāḥa riyāḍa ṣiḥḥiyya.}
\end{quote}

\begin{quote}
Swimming is a healthy sport.
\end{quote}

\begin{cousinweb}
\textbf{\ar{سباحة}} (\emph{sibāḥa}) is built on the native Arabic root \emph{s-b-ḥ}, \textquotedblleft{}to swim, to float\textquotedblright{} (Lane). Sabaḥa, to swim, which you know, is from the same root.
\end{cousinweb}

\begin{grammarlens}[title={What you've built}]
One more word for describing what you lived and felt.
\end{grammarlens}

\subsection*{Your turn}
\begin{itemize}
\raggedright
  \item \emph{Say it:} \emph{sibāḥa}
  \item \emph{Say it:} \emph{sibāḥa}, once more
  \item \emph{Say it:} \emph{as-sibāḥa riyāḍa ṣiḥḥiyya.} — Swimming is a healthy sport.
  \item \emph{Recall:} say the Arabic for a medal, then the Arabic for to qualify, then say \emph{sibāḥa} again
\end{itemize}

\begin{tcolorbox}[breakable,colback=teal!4,colframe=teal!35!black,title={Before you move on}]
What does \ar{سباحة} mean? (\textquotedblleft{}Swimming\textquotedblright{}.) Say it once more.
\end{tcolorbox}
\section[\texorpdfstring{\ar{ملاكمة} (mulākama)}{ملاكمة (mulākama)}]{\ar{ملاكمة} (mulākama) — boxing}

\label{lesson:AR-C548-mulakama}



\begin{quote}
Before the new one: say the Arabic for to qualify, then the Arabic for swimming.
\end{quote}

\subsection*{You'll want to know: \ar{ملاكمة}}
\textbf{\ar{ملاكمة}} — \emph{mulākama} — \textquotedblleft{}boxing\textquotedblright{}.

\begin{quote}
\ar{يشاهد} \ar{أبي} \ar{الملاكمة} \ar{على} \ar{التلفاز}.
\end{quote}

\begin{quote}
\emph{yushāhidu abī al-mulākama ʿalā at-tilfāz.}
\end{quote}

\begin{quote}
My father watches boxing on television.
\end{quote}

\begin{cousinweb}
\textbf{\ar{ملاكمة}} (\emph{mulākama}) is built on the native Arabic root \emph{l-k-m}, \textquotedblleft{}to strike with the fist\textquotedblright{} (Lane). Mulākama is the verbal noun of Form III: hitting one another with the fist.
\end{cousinweb}

\begin{grammarlens}[title={What you've built}]
One more word for describing what you lived and felt.
\end{grammarlens}

\subsection*{Your turn}
\begin{itemize}
\raggedright
  \item \emph{Say it:} \emph{mulākama}
  \item \emph{Say it:} \emph{mulākama}, once more
  \item \emph{Say it:} \emph{yushāhidu abī al-mulākama ʿalā at-tilfāz.} — My father watches boxing on television.
  \item \emph{Recall:} say the Arabic for to qualify, then the Arabic for swimming, then say \emph{mulākama} again
\end{itemize}

\begin{tcolorbox}[breakable,colback=teal!4,colframe=teal!35!black,title={Before you move on}]
What does \ar{ملاكمة} mean? (\textquotedblleft{}Boxing\textquotedblright{}.) Say it once more.
\end{tcolorbox}
\section[\texorpdfstring{\ar{صارع} (ṣāraʿa)}{صارع (ṣāraʿa)}]{\ar{صارع} (ṣāraʿa) — to wrestle (with)}

\label{lesson:AR-C548-saraa}



\begin{quote}
Before the new one: say the Arabic for swimming, then the Arabic for boxing.
\end{quote}

\subsection*{You'll want to know: \ar{صارع}}
\textbf{\ar{صارع}} — \emph{ṣāraʿa} — \textquotedblleft{}to wrestle (with)\textquotedblright{}.

\begin{quote}
\ar{صارع} \ar{اللاعب} \ar{بطل} \ar{البلد}.
\end{quote}

\begin{quote}
\emph{ṣāraʿa al-lāʿib baṭal al-balad.}
\end{quote}

\begin{quote}
The player wrestled the national champion.
\end{quote}

Present: \ar{يصارع}, \emph{yuṣāriʿ}.

\begin{cousinweb}
\textbf{\ar{صارع}} (\emph{ṣāraʿa}) is built on the native Arabic root \emph{ṣ-r-ʿ}, \textquotedblleft{}to throw down\textquotedblright{} (Lane). Form III ṣāraʿa is to try to throw one another down.
\end{cousinweb}

\begin{grammarlens}[title={What you've built}]
One more word for describing what you lived and felt.
\end{grammarlens}

\subsection*{Your turn}
\begin{itemize}
\raggedright
  \item \emph{Say it:} \emph{ṣāraʿa}
  \item \emph{Say it:} \emph{ṣāraʿa}, once more
  \item \emph{Say it:} \emph{ṣāraʿa al-lāʿib baṭal al-balad.} — The player wrestled the national champion.
  \item \emph{Recall:} say the Arabic for swimming, then the Arabic for boxing, then say \emph{ṣāraʿa} again
\end{itemize}

\begin{tcolorbox}[breakable,colback=teal!4,colframe=teal!35!black,title={Before you move on}]
What does \ar{صارع} mean? (\textquotedblleft{}To wrestle (with)\textquotedblright{}.) Say it once more.
\end{tcolorbox}
\section[\texorpdfstring{\ar{نافس} (nāfasa)}{نافس (nāfasa)}]{\ar{نافس} (nāfasa) — to compete with}

\label{lesson:AR-C548-nafasa}



\begin{quote}
Before the new one: say the Arabic for boxing, then the Arabic for to wrestle.
\end{quote}

\subsection*{You'll want to know: \ar{نافس}}
\textbf{\ar{نافس}} — \emph{nāfasa} — \textquotedblleft{}to compete with\textquotedblright{}.

\begin{quote}
\ar{نافست} \ar{أختي} \ar{عشر} \ar{بنات} \ar{في} \ar{السباق}.
\end{quote}

\begin{quote}
\emph{nāfasat ukhtī ʿashr banāt fī as-sibāq.}
\end{quote}

\begin{quote}
My sister competed with ten girls in the race.
\end{quote}

Present: \ar{ينافس}, \emph{yunāfis}.

\begin{cousinweb}
\textbf{\ar{نافس}} (\emph{nāfasa}) is built on the native Arabic root \emph{n-f-s}, \textquotedblleft{}a soul; to be precious\textquotedblright{} (Lane). Form III nāfasa is to vie with someone for a precious thing.
\end{cousinweb}

\begin{grammarlens}[title={What you've built}]
One more word for describing what you lived and felt.
\end{grammarlens}

\subsection*{Your turn}
\begin{itemize}
\raggedright
  \item \emph{Say it:} \emph{nāfasa}
  \item \emph{Say it:} \emph{nāfasa}, once more
  \item \emph{Say it:} \emph{nāfasat ukhtī ʿashr banāt fī as-sibāq.} — My sister competed with ten girls in the race.
  \item \emph{Recall:} say the Arabic for boxing, then the Arabic for to wrestle, then say \emph{nāfasa} again
\end{itemize}

\begin{tcolorbox}[breakable,colback=teal!4,colframe=teal!35!black,title={Before you move on}]
What does \ar{نافس} mean? (\textquotedblleft{}To compete with\textquotedblright{}.) Say it once more.
\end{tcolorbox}
\section[\texorpdfstring{\ar{لكم} (lakama)}{لكم (lakama)}]{\ar{لكم} (lakama) — to punch}

\label{lesson:AR-C548-lakama}



\begin{quote}
Before the new one: say the Arabic for to wrestle, then the Arabic for to compete with.
\end{quote}

\subsection*{You'll want to know: \ar{لكم}}
\textbf{\ar{لكم}} — \emph{lakama} — \textquotedblleft{}to punch\textquotedblright{}.

\begin{quote}
\ar{لكم} \ar{الملاكم} \ar{خصمه}.
\end{quote}

\begin{quote}
\emph{lakama al-mulākim khaṣmahu.}
\end{quote}

\begin{quote}
The boxer punched his opponent.
\end{quote}

Present: \ar{يلكم}, \emph{yalkum}.

\begin{cousinweb}
\textbf{\ar{لكم}} (\emph{lakama}) is built on the native Arabic root \emph{l-k-m}, \textquotedblleft{}to strike with the fist\textquotedblright{} (Lane). Lane gives lakma as a blow with the fist; mulākama, boxing, is from the same root.
\end{cousinweb}

\begin{grammarlens}[title={What you've built}]
One more word for describing what you lived and felt.
\end{grammarlens}

\subsection*{Your turn}
\begin{itemize}
\raggedright
  \item \emph{Say it:} \emph{lakama}
  \item \emph{Say it:} \emph{lakama}, once more
  \item \emph{Say it:} \emph{lakama al-mulākim khaṣmahu.} — The boxer punched his opponent.
  \item \emph{Recall:} say the Arabic for to wrestle, then the Arabic for to compete with, then say \emph{lakama} again
\end{itemize}

\begin{tcolorbox}[breakable,colback=teal!4,colframe=teal!35!black,title={Before you move on}]
What does \ar{لكم} mean? (\textquotedblleft{}To punch\textquotedblright{}.) Say it once more.
\end{tcolorbox}
